%=====================================================================
\chapter{The Virtual Machine Monitor}\label{ch:vmm}
%=====================================================================
\locator{1}{Part I --- Foundations of Virtualization}

\begin{outcomes}
By the end of this chapter you will be able to:
\begin{tight}
  \item \textbf{State} the three properties Popek and Goldberg require of a
        virtual machine monitor, and \textbf{explain} what each one rules out.
  \item \textbf{Distinguish} privileged from sensitive instructions, and
        \textbf{state} the theorem that connects them to virtualizability.
  \item \textbf{Explain} why x86 failed the theorem's condition, using
        \texttt{POPF} as the worked case, and \textbf{say} precisely what goes
        wrong --- which is not a crash.
  \item \textbf{Describe} binary translation and ring deprivileging as the two
        software answers, and \textbf{compare} what each costs.
  \item \textbf{Apply} the theorem to an architecture you are given a
        specification for, and decide whether trap-and-emulate alone would
        suffice.
\end{tight}
\end{outcomes}

\begin{prereq}
\chref{what} (the layer and what it promises) and \chref{architectures} (the
three techniques, named but not yet justified). You also need, from Operating
Systems, what a privilege level is and what happens on a trap. If that is
shaky, read the first page of Section~3.1 slowly; everything after it depends
on the idea that the processor behaves differently according to who is asking.
\end{prereq}

\begin{keyterms}
virtual machine monitor $\bullet$ equivalence $\bullet$ resource control
$\bullet$ efficiency $\bullet$ privileged instruction $\bullet$ sensitive
instruction $\bullet$ control-sensitive $\bullet$ behaviour-sensitive $\bullet$
trap $\bullet$ trap-and-emulate $\bullet$ Popek--Goldberg theorem $\bullet$
silent failure $\bullet$ ring deprivileging $\bullet$ ring aliasing $\bullet$
binary translation $\bullet$ translation cache $\bullet$ hybrid virtualization
\end{keyterms}

\section{What a Monitor Must Guarantee}

In 1974 Gerald Popek and Robert Goldberg asked a question that had not been
asked precisely before: \emph{what has to be true of a processor for it to be
possible to write a correct virtual machine monitor for it?} Their answer is
the only theorem in this book, and it explains the twenty-five-year gap in
\dsref{timeline}.

They begin by saying what a monitor is \emph{for}. Three properties, all three
required.

\begin{definitionbox}[The Three Properties of a Virtual Machine Monitor]
\begin{tight}
  \item \term{Equivalence}. A program running under the monitor behaves
        identically to the same program running directly on the hardware,
        except for timing and for resources that are genuinely absent. This is
        the ``preserved interface'' of \chref{what}, stated formally.
  \item \term{Resource control}. The monitor is in complete control of the
        resources it virtualizes. A guest cannot reach memory it was not
        given, cannot affect another guest, and cannot take control away from
        the monitor.
  \item \term{Efficiency}. A \emph{statistically dominant} subset of the
        guest's instructions executes directly on the real processor, with no
        intervention.
\end{tight}
\end{definitionbox}

\begin{conceptbox}[Efficiency Is the Clause That Does the Work]
Drop the third property and the problem becomes trivial: write an interpreter.
An emulator satisfies equivalence and resource control perfectly, because it
never lets the guest touch the processor at all. It is also fifty times slower,
which is why nobody consolidates a data centre with one.

\smallskip
So the entire difficulty of this chapter lives in the word \emph{dominant}. The
monitor must let the guest run on the real silicon almost all of the time, and
must nevertheless remain in control --- and those two demands pull in opposite
directions. The rest of the chapter is about where exactly they tear.

\smallskip
Notice also what equivalence \emph{excludes}: timing. A guest may legitimately
observe that it is running slower, and a guest that reads the cycle counter can
often tell it is virtualized. Popek and Goldberg were explicit that this is
allowed. Software that uses timing to \emph{detect} a hypervisor --- malware
avoiding analysis, licence enforcement --- is exploiting a gap the theorem
deliberately left open.
\end{conceptbox}

\section{Privileged, Sensitive, and the Difference}

The classical way to build a monitor is \term{trap-and-emulate}: run the guest
at a reduced privilege level, let the hardware trap anything it may not do, and
have the monitor emulate the trapped instruction in a way consistent with the
guest's virtual machine.

\dsfig{%
\begin{tikzpicture}[font=\scriptsize,
  st/.style={rounded corners=2pt, line width=0.8pt, align=center,
             minimum width=3.9cm, minimum height=0.64cm, inner sep=3pt},
  g/.style={st, draw=secondary, fill=secondary!8, text=secondary!80!black},
  h/.style={st, draw=primary, fill=primary!12, text=primary},
  p/.style={st, draw=accent, fill=accent!8, text=accent!70!black},
  ar/.style={-{Stealth[length=2mm]}, line width=0.9pt},
  note/.style={font=\scriptsize\itshape, text=neutral, anchor=west, align=left}
]
  \node[g] (g1) at (0,4.0)  {guest kernel executes\\an instruction};
  \node[p] (t1) at (0,2.85) {may it? the \emph{processor} decides};
  \node[g] (ok) at (-3.1,1.6) {ordinary instruction:\\runs at full speed};
  \node[h] (tr) at (3.1,1.6)  {privileged: \textbf{trap}};
  \node[h] (em) at (3.1,0.45) {monitor emulates it\\against the guest's\\virtual state};
  \node[g] (rt) at (3.1,-0.9) {control returns to\\the guest};

  \draw[ar, secondary!60] (g1) -- (t1);
  \draw[ar, greenhl!70] (t1) -- (ok);
  \draw[ar, accent!70]  (t1) -- (tr);
  \draw[ar, primary!60] (tr) -- (em);
  \draw[ar, primary!60] (em) -- (rt);
  \draw[ar, greenhl!70] (ok) to[out=-90,in=180] (0,-1.6) -- (1.2,-1.6);
  \draw[ar, primary!60] (rt) to[out=-90,in=0] (0,-1.9) -- (-1.2,-1.9);

  \node[note] at (5.4,1.6)  {the \emph{dominant} case\\--- efficiency};
  \node[note] at (5.4,0.45) {the \emph{rare} case\\--- equivalence and control};
  \node[note] at (-6.4,2.85) {this is the step\\x86 could not do};
  \draw[accent!60, line width=0.8pt, -{Stealth[length=1.6mm]}]
        (-3.5,2.85) -- (-2.2,2.85);
\end{tikzpicture}}%
{Trap-and-emulate. The design is sound and was old by 1974; the whole question
is whether the processor traps everything the monitor needs it to
trap.}{trap-emulate}

For this to work, the processor must trap \emph{every} instruction whose
behaviour the monitor needs to control. Popek and Goldberg made that precise
with two definitions.

\begin{definitionbox}[Privileged and Sensitive Instructions]
An instruction is \term{privileged} if it traps when executed in user mode and
does not trap in supervisor mode. This is a property of the \emph{hardware}:
the processor is built to refuse it.

\smallskip
An instruction is \term{sensitive} if it is one the monitor must know about.
There are two kinds:
\begin{tight}
  \item \term{control-sensitive} --- it changes the configuration of resources:
        the page tables, the interrupt mask, the privilege level.
  \item \term{behaviour-sensitive} --- its \emph{result} depends on the
        configuration: it reads a register whose value differs between
        privilege levels, or depends on where in memory it is running.
\end{tight}

\smallskip
Sensitivity is a property of what the \emph{monitor needs}; privilege is a
property of what the \emph{hardware does}. Nothing guarantees that they
coincide.
\end{definitionbox}

\begin{conceptbox}[The Theorem]
\textbf{Popek and Goldberg, 1974.} For any conventional third-generation
computer, a virtual machine monitor may be constructed if the set of sensitive
instructions is a \emph{subset} of the set of privileged instructions.

\smallskip
That is the whole condition. Every instruction the monitor needs to know about
must be one the hardware will trap. If so, trap-and-emulate works: run the guest
unprivileged, catch the sensitive instructions as traps, emulate them, and
everything else runs at full speed --- equivalence, control and efficiency
together.

\smallskip
And the contrapositive is the interesting half. If even one instruction is
sensitive but \emph{not} privileged, the guest can execute it, the hardware will
not trap it, and the monitor will never find out.
\end{conceptbox}

\dsfig{%
\begin{tikzpicture}[font=\small]
  % --- virtualizable ---
  \node[font=\small\bfseries, text=greenhl!55!black] at (0,3.15)
       {Virtualizable: sensitive $\subseteq$ privileged};
  \draw[draw=primary, line width=1.1pt, fill=primary!8]
        (-2.6,0.1) rectangle (2.6,2.6);
  \node[text=primary, font=\scriptsize\bfseries, anchor=north west]
       at (-2.5,2.52) {privileged};
  \draw[draw=accent, line width=1.1pt, fill=accent!10]
        (-1.5,0.5) rectangle (1.5,1.95);
  \node[text=accent!70!black, font=\scriptsize\bfseries] at (0,1.55) {sensitive};
  \node[font=\scriptsize\itshape, text=greenhl!55!black] at (0,1.0)
       {every one of these traps};
  \node[font=\scriptsize\itshape, text=neutral, anchor=north] at (0,-0.1)
       {trap-and-emulate alone is enough};

  % --- not virtualizable ---
  \node[font=\small\bfseries, text=accent!70!black] at (8.6,3.15)
       {x86 before 2005: 17 instructions outside};
  \draw[draw=primary, line width=1.1pt, fill=primary!8]
        (6.0,0.1) rectangle (11.2,2.6);
  \node[text=primary, font=\scriptsize\bfseries, anchor=north west]
       at (6.1,2.52) {privileged};
  \draw[draw=accent, line width=1.1pt, fill=accent!10]
        (7.4,0.5) rectangle (12.6,1.95);
  \node[text=accent!70!black, font=\scriptsize\bfseries] at (9.0,1.55) {sensitive};
  \node[font=\scriptsize\bfseries, text=accent, align=center] at (11.9,1.22)
       {\texttt{POPF}\\\texttt{SGDT}\\\texttt{PUSHF}\\\ldots};
  \draw[decorate, decoration={brace, amplitude=4pt}, accent, line width=0.9pt]
        (12.7,1.95) -- (12.7,0.5);
  \node[font=\scriptsize\itshape, text=accent!70!black, anchor=west, align=left]
       at (12.95,1.22) {sensitive but\\not privileged:\\no trap, no\\notification};
  \node[font=\scriptsize\itshape, text=neutral, anchor=north] at (8.6,-0.1)
       {trap-and-emulate silently produces wrong answers};
\end{tikzpicture}}%
{The condition, and the architecture that failed it. The seventeen instructions
on the right are not exotic: several are ordinary, unprivileged instructions a
compiler emits without comment.}{popek-goldberg}

\section{Why x86 Failed}

In 2000 John Robin and Cynthia Irvine went through the Pentium instruction set
against Popek and Goldberg's condition and found \textbf{seventeen}
instructions that are sensitive but not privileged. The IBM System/370 had
satisfied the condition; x86, designed for a single operating system on a
desktop, had never needed to.

The clearest of the seventeen is \texttt{POPF}.

\begin{worked}[How \texttt{POPF} Breaks a Hypervisor Without Crashing It]
\texttt{POPF} pops a word off the stack into the \texttt{EFLAGS} register.
\texttt{EFLAGS} holds ordinary arithmetic flags --- carry, zero, sign --- and
also the \term{interrupt flag} \texttt{IF}, which controls whether the processor
accepts hardware interrupts.

\smallskip
\textbf{What the guest kernel intends.} It is entering a critical section and
wants interrupts off:
\begin{tight}
  \item push a flags word with \texttt{IF = 0};
  \item execute \texttt{POPF};
  \item expect interrupts to be disabled.
\end{tight}

\smallskip
\textbf{Step 1 --- is \texttt{POPF} sensitive?} Yes, control-sensitive: it
changes the interrupt configuration, which the monitor must control. If a guest
could really disable interrupts, it could refuse to give the processor back,
destroying resource control.

\smallskip
\textbf{Step 2 --- is \texttt{POPF} privileged?} \textbf{No.} In user mode the
x86 \texttt{POPF} does not fault. It updates the arithmetic flags and
\emph{silently ignores the write to} \texttt{IF}. The designers chose this
deliberately in 1985, so that user programs manipulating arithmetic flags would
not need to care about \texttt{IF}.

\smallskip
\textbf{Step 3 --- so what happens under a monitor?} The guest kernel runs
unprivileged. It executes \texttt{POPF}:
\begin{tight}
  \item no trap occurs --- the monitor is never told;
  \item the arithmetic flags are updated, so the instruction appears to work;
  \item \texttt{IF} is unchanged, so interrupts remain \emph{enabled};
  \item the instruction completes and the guest proceeds, believing it is in a
        critical section.
\end{tight}

\smallskip
\textbf{Step 4 --- what the failure looks like.} Not a crash. An interrupt
arrives during the critical section the guest believes it has protected. A data
structure is modified halfway through an update. Perhaps nothing happens. Perhaps
the guest's scheduler corrupts its run queue, at a rate of once a fortnight, in
a way that is not reproducible and cannot be traced to the instruction that
caused it.

\smallskip
\textbf{The point.} A \emph{crash} would be a gift --- it is a signal. The
seventeen instructions do something far worse: they make the equivalence
property fail \emph{quietly}. The guest cannot detect it, the monitor cannot
detect it, and the resulting bug reports are unfalsifiable. This, and not
performance, is why x86 was called non-virtualizable for twenty-five years.
\end{worked}

\begin{alertbox}[``Non-virtualizable'' never meant impossible]
It meant that \emph{trap-and-emulate alone} is not sufficient --- that no
correct monitor can be built the classical way. VMware shipped a correct x86
hypervisor in 1999 and Xen in 2003, so clearly monitors could be built. What
the theorem denies is that they can be built cheaply and simply. Both products
are, in effect, elaborate workarounds for seventeen instructions, and both
workarounds cost something the next section quantifies.
\end{alertbox}

\section{Two Ways Round}

\subsection{Binary translation}

VMware's answer, and the one that founded the industry. The monitor never lets
the guest kernel's instructions reach the processor unexamined. Instead it
reads the kernel's code a basic block at a time, \emph{translates} it ---
copying ordinary instructions unchanged and replacing each of the seventeen
with a call into the monitor --- and executes the translated block. Translations
are cached, so a hot loop is translated once and then runs many times.

Guest \emph{user} code is not translated at all. It runs directly, at full
speed, because an application executing \texttt{POPF} and losing the write to
\texttt{IF} is exactly what would happen on real hardware: equivalence is
preserved for free.

\begin{conceptbox}[Why Translation Was Fast Enough to Sell]
It sounds ruinous and it was not, for two reasons.

\smallskip
First, the ratio. A kernel is a small fraction of what runs on a busy machine,
and the translation cost is paid once per block rather than once per execution.

\smallskip
Second --- and this surprised people --- a translator can be \emph{faster} than
the hardware trap it replaces. A trap on 2005-era x86 cost well over a thousand
cycles of saving and restoring state. A translated call into the monitor cost
tens. For several years after hardware assistance arrived, VMware's translator
beat the first generation of VT-x on system-call-heavy workloads, and VMware
said so publicly. \chref{hardware} is the story of Intel and AMD closing that
gap.
\end{conceptbox}

\subsection{Ring deprivileging and paravirtualization}

Xen's answer. x86 provides four privilege levels --- rings 0 to 3 --- of which
operating systems historically used only two. Xen put the hypervisor in ring~0,
moved the guest kernel to ring~1, and left applications in ring~3.

\dsfig{%
\begin{tikzpicture}[font=\scriptsize,
  ring/.style={rounded corners=2pt, line width=0.9pt, align=center,
               minimum width=3.5cm, minimum height=0.7cm, inner sep=3pt},
  lbl/.style={font=\scriptsize\bfseries, text=neutral!40!black, anchor=east},
  note/.style={font=\scriptsize\itshape, text=accent, anchor=west, align=left}
]
  % native
  \node[font=\small\bfseries, text=primary] at (2.3,3.0) {Native};
  \node[lbl] at (0.2,2.2) {ring 0}; \node[lbl] at (0.2,1.3) {ring 1};
  \node[lbl] at (0.2,0.4) {ring 2}; \node[lbl] at (0.2,-0.5) {ring 3};
  \node[ring, draw=purplehl, fill=purplehl!12, text=purplehl!70!black]
       at (2.3,2.2) {OS kernel};
  \node[ring, draw=neutral!50, fill=neutral!6, text=neutral!45!black, dashed]
       at (2.3,1.3) {unused};
  \node[ring, draw=neutral!50, fill=neutral!6, text=neutral!45!black, dashed]
       at (2.3,0.4) {unused};
  \node[ring, draw=secondary, fill=secondary!8, text=secondary!80!black]
       at (2.3,-0.5) {applications};

  % deprivileged
  \node[font=\small\bfseries, text=primary] at (9.0,3.0) {Xen, ring deprivileged};
  \node[ring, draw=primary, fill=primary!14, text=primary]
       at (9.0,2.2) {\textbf{hypervisor}};
  \node[ring, draw=purplehl, fill=purplehl!12, text=purplehl!70!black]
       at (9.0,1.3) {guest kernel};
  \node[ring, draw=neutral!50, fill=neutral!6, text=neutral!45!black, dashed]
       at (9.0,0.4) {unused};
  \node[ring, draw=secondary, fill=secondary!8, text=secondary!80!black]
       at (9.0,-0.5) {applications};

  \draw[-{Stealth[length=2.2mm]}, accent, line width=1pt]
        (4.3,2.2) to[out=20,in=160] (7.1,1.3);
  \node[note] at (11.0,1.3) {now unprivileged,\\so most sensitive\\instructions
                             do trap};
  \node[note] at (11.0,2.2) {and the ones that\\still do not are\\replaced by
                             hypercalls};
\end{tikzpicture}}%
{Ring deprivileging. Moving the guest kernel down one ring makes the hardware
trap most of what the monitor needs; the instructions that still slip through
are removed from the guest's source instead.}{ring-deprivileging}

Deprivileging alone does not solve the problem --- \texttt{POPF} still fails
silently in ring~1 --- so Xen took the second step: it \emph{modified the guest
kernel}. Where Linux would have executed one of the seventeen, the Xen port
issues a \term{hypercall} instead. That is paravirtualization, and the price is
exactly the one named in \chref{architectures}: you need the guest's source.
Xen ran Linux and NetBSD from the first release and could not run Windows at
all.

\begin{alertbox}[Ring aliasing: the guest can tell]
Deprivileging breaks equivalence in a second, subtler way. The unprivileged
instruction \texttt{PUSHF} pushes \texttt{EFLAGS} --- including the current
privilege level --- onto the stack, and \texttt{SGDT} and \texttt{SIDT} read
descriptor-table registers. None of them traps. A guest kernel that executes
any of them sees ring~1 where it expects ring~0, and a descriptor table that is
not its own.

\smallskip
This is \term{ring aliasing}, and it is why \texttt{SGDT} is on the list of
seventeen. It also makes hypervisor detection trivial: a few unprivileged
instructions and the guest knows. \chref{confidential} is where that question
comes back inverted --- there, the guest \emph{wants} to prove what it is
running on.
\end{alertbox}

\rowcolors{2}{white}{lightbg}
\begin{longtable}{@{}L{3.3cm}L{3.8cm}L{3.8cm}L{3.8cm}@{}}
\toprule
\rowcolor{primary!15}
& \textbf{Binary translation} & \textbf{Paravirtualization} &
\textbf{Hardware assistance} \\
\midrule
\endhead
Who changes & The monitor, at run time & The guest, at compile time & The
processor, at design time \\
Unmodified guests & Yes & \textbf{No} & Yes \\
Where the cost falls & Translating, and a cache to maintain & Nothing at run
time; everything at porting time & A trap per sensitive instruction \\
Monitor complexity & Very high --- a working x86 translator & Low & Low \\
Shipped & VMware, 1999 & Xen, 2003 & Intel VT-x 2005, AMD-V 2006 \\
Still used & Where hardware assistance is absent or slower & As \emph{drivers},
almost universally (\chref{paravirt}) & Everywhere (\chref{hardware}) \\
\bottomrule
\end{longtable}

\begin{pitfall}
\begin{tight}
  \item \textbf{Saying that non-virtualizable means ``cannot be
        virtualized''.} It means trap-and-emulate alone is insufficient. Two
        companies built correct monitors anyway, at considerable cost.
  \item \textbf{Expecting the failure to be a crash.} The \texttt{POPF} failure
        is silent and intermittent. Students who look for an exception in the
        worked example have missed what makes those seventeen instructions
        dangerous.
  \item \textbf{Confusing privileged with sensitive.} Privileged is what the
        hardware traps; sensitive is what the monitor needs. The theorem is
        precisely the claim that the first set must contain the second.
  \item \textbf{Believing hardware assistance made the theorem obsolete.} It
        made x86 satisfy the condition. The theorem is what tells you
        \emph{why} VT-x was needed and what it had to do --- and it still
        decides the question for any new architecture, which is why RISC-V
        designers argued about exactly this.
  \item \textbf{Assuming efficiency means ``fast''.} It is a formal property:
        a statistically dominant subset of instructions runs natively. An
        emulator fails it at any speed.
\end{tight}
\end{pitfall}

\begin{lab}[Catching a Sensitive Instruction in the Act]
No hypervisor needed: this runs on the host and shows the hardware behaviour
the whole chapter rests on. Needs \texttt{gcc}. Expect twenty minutes.

\begin{lstlisting}[language=C]
/* ch03_popf.c -- show that POPF fails silently in user mode.
   Build: gcc -O0 -o ch03_popf ch03_popf.c     Run: ./ch03_popf   */
#include <stdio.h>

/* --- 1. Read EFLAGS, so we can see what the processor really did. */
static unsigned long read_flags(void) {
    unsigned long f;
    __asm__ volatile ("pushfq ; popq %0" : "=r"(f));
    return f;
}

/* --- 2. Try to write EFLAGS, clearing the interrupt flag (bit 9). */
static void write_flags(unsigned long f) {
    __asm__ volatile ("pushq %0 ; popfq" : : "r"(f) : "cc");
}

int main(void) {
    unsigned long before = read_flags();
    printf("before : EFLAGS=0x%08lx  IF=%lu\n", before, (before >> 9) & 1);

    /* --- 3. Ask for interrupts off, and for the carry flag on.
       A kernel entering a critical section does exactly this. */
    unsigned long want = (before & ~(1UL << 9)) | 1UL;   /* IF=0, CF=1 */
    write_flags(want);

    /* --- 4. Look at what we actually got. */
    unsigned long after = read_flags();
    printf("wanted : EFLAGS=0x%08lx  IF=%lu  CF=%lu\n",
           want, (want >> 9) & 1, want & 1);
    printf("after  : EFLAGS=0x%08lx  IF=%lu  CF=%lu\n",
           after, (after >> 9) & 1, after & 1);

    /* --- 5. The verdict. No signal was raised; nothing faulted. */
    printf("\ncarry flag taken?     %s\n",
           (after & 1) ? "yes -- the instruction 'worked'" : "no");
    printf("interrupt flag taken? %s\n",
           (((after >> 9) & 1) == 0) ? "yes" :
           "NO -- written, ignored, and no trap. This is the bug.");
    return 0;
}
\end{lstlisting}

\textbf{What to notice.} The program does not crash and reports no error. The
carry flag changed, so by every signal available to the program the instruction
succeeded --- and the interrupt flag did not change. Now reread Step~4 of the
worked example: a guest kernel in this position believes it holds a lock it does
not hold.
\end{lab}

\begin{checkpoint}
\begin{tightnum}
  \item An architecture has an instruction \texttt{RDCFG} that returns the
        current privilege level and does not fault in user mode. Classify it
        using the two definitions, and say whether this architecture satisfies
        the Popek--Goldberg condition.
  \item Why would a \emph{crash} on executing \texttt{POPF} in a guest have
        been a better outcome for hypervisor authors than what x86 actually
        does?
  \item Binary translation needs no cooperation from the guest, and
        paravirtualization needs the guest's source code. Give one situation
        where the second is nevertheless the better engineering choice, and say
        why.
\end{tightnum}
\end{checkpoint}

\begin{chaptersummary}
\begin{tight}
  \item A virtual machine monitor must provide equivalence, resource control
        and efficiency. Efficiency --- a statistically dominant subset running
        natively --- is what rules out an emulator and makes the problem hard.
  \item Privileged means the hardware traps it; sensitive means the monitor
        needs to know about it. They are different properties of different
        things.
  \item Popek and Goldberg: a monitor can be built by trap-and-emulate if the
        sensitive instructions are a subset of the privileged ones.
  \item x86 had seventeen instructions that are sensitive and not privileged.
        \texttt{POPF} silently discards a write to the interrupt flag in user
        mode, so the guest believes it disabled interrupts and did not.
  \item The failure is silent, intermittent and untraceable --- which is worse
        than a crash, and is why x86 was called non-virtualizable.
  \item Binary translation rewrites the guest's kernel code at run time, needs
        no guest changes, and was for some years faster than the first
        hardware assistance. Ring deprivileging plus paravirtualization is
        simpler and needs the guest's source.
  \item The theorem is not obsolete. It is what \chref{hardware} is an answer
        to, and it still decides the question for any new architecture.
\end{tight}
\end{chaptersummary}

\begin{reviewq}
  \item \textbf{(MCQ)} Popek and Goldberg's efficiency property requires that:
        (a)~the monitor adds less than 5\% overhead; (b)~a statistically
        dominant subset of guest instructions executes directly on the
        processor; (c)~the guest boots in under a minute; (d)~no instruction is
        ever trapped.
  \item \textbf{(MCQ)} An instruction that is sensitive but not privileged
        causes: (a)~a general protection fault; (b)~the guest to crash;
        (c)~the monitor to lose control or correctness with no notification;
        (d)~a performance penalty only.
  \item \textbf{(Short)} State the Popek--Goldberg condition in one sentence,
        and explain what each of the two instruction sets is a property of.
  \item \textbf{(Short)} Describe what \texttt{POPF} does in user mode on x86
        and why that single design decision cost the industry twenty-five
        years.
  \item \textbf{(Short)} Explain how binary translation could outperform
        hardware-assisted virtualization on a system-call-heavy workload in
        2006, even though the hardware was doing the same job.
  \item \textbf{(Applied)} You are reviewing the specification of a new
        64-bit architecture. Write the procedure you would follow to decide
        whether it can be virtualized by trap-and-emulate alone, say what
        evidence you would need from the vendor, and state what you would
        recommend if you found three offending instructions.
  \item \textbf{(Applied)} A colleague proposes to run an unmodified legacy
        operating system, whose vendor no longer exists, on current hardware.
        Using this chapter and \chref{architectures}, set out which
        architectures are available to them, which are ruled out and why, and
        what you would check first.
\end{reviewq}
